\section{Результаты диагностики и найденные недочёты}
\TemplateNote{Числа и результаты проверок пока не внесены.
Заполнять раздел только по выполненным запускам и сохранённым логам.}

\begin{tabularx}{\linewidth}{@{}p{34mm}YY@{}}
\hline
Проверка & Результат & Подтверждение \\
\hline
Инварианты дерева & \Field{заполнить} & \Field{тест, seed, лог} \\
Сравнение с эталоном & \Field{заполнить} & \Field{тест, seed, лог} \\
Ошибки памяти & \Field{заполнить} & \Field{команда и лог} \\
Save и Load & \Field{заполнить} & \Field{файлы и вывод} \\
Ошибочный Load & \Field{заполнить} & \Field{проверка прежних данных} \\
\hline
\end{tabularx}

\subsection{Карточка обнаруженной проблемы}
\textbf{Проявление:} \Field{неправильный ответ, диагностика или измеренный недостаток}.

\textbf{Воспроизведение:} \Field{коммит, вход, seed, команда}.

\textbf{Причина:} \Field{участок и нарушенное свойство}.

\textbf{Исправление:} \Field{что изменено и почему}.

\textbf{Повторная проверка:} \Field{тот же тест, результат и лог}.

Карточку повторить для каждой реальной проблемы. Если проблем не найдено,
описать проверки и их границы; вымышленные дефекты не добавлять.

\section{Сравнение версий по времени и памяти}
\begin{tabularx}{\linewidth}{@{}p{29mm}YYY@{}}
\hline
Показатель & Базовая версия & Исправленная & Нагрузка и лог \\
\hline
Время, медиана & \Field{с} & \Field{с} & \Field{ссылка} \\
Разброс времени & \Field{с} & \Field{с} & \Field{ссылка} \\
Пик RSS & \Field{единицы} & \Field{единицы} & \Field{ссылка} \\
Пик кучи & \Field{единицы} & \Field{единицы} & \Field{ссылка} \\
Ошибки памяти & \Field{результат} & \Field{результат} & \Field{ссылка} \\
\hline
\end{tabularx}

Ускорение при сопоставимых условиях: $T_{\text{до}}/T_{\text{после}}$.
Изменение времени в процентах:
$(T_{\text{до}}-T_{\text{после}})/T_{\text{до}}\cdot100\%$.
Изменение памяти анализировать отдельно.

\Field{Добавить сценарии и размеры данных, при необходимости графики
и сравнение степеней. Объяснить причины различий.
Если базовая версия некорректна, отметить это при сравнении.}
